\section{Methods}
\label{sec:methods}

\subsection{Research design}

The study is a pre-specified, factorial, multi-country benchmark experiment
quantifying geographic and persona-modulated bias in 14 large language models on
air pollution policy tasks. The pre-specified design crosses 15 countries
(stratified on three theoretically grounded axes), 14 LLMs spanning five access
tiers, one domain, five task types and two persona conditions, with two
stochastic replications per cell; a post-hoc extension
(Section~\ref{sec:expansion}) enlarges the sample to 25 countries. The stimulus
set crosses each country's five tasks and two personas in English, plus a
within-country native-language rendering for the nine countries with a target
native language, scored on the $[0,1]$ composite. The protocol fixed, before
collection, the hypotheses, primary models, persona protocol, exclusion criteria
and multiple-comparison corrections. A Brazil-only collection executed under an
earlier three-domain design is reclassified as an extended pilot and reported
separately; it informed prompt calibration and contributes no observations to
the main analysis.

\subsection{Country selection}

Countries were selected by theoretical sampling~\citep{patton2015qualitative}
triangulated across three independent classifications: the Global North/South
distinction, operationalised through UNCTAD's developed/developing economies
classification~\citep{unctad2024classification}, which carries the
coloniality-of-knowledge framework motivating the
study~\citep{mohamed2020decolonial}; the Joshi et al.\ six-class taxonomy of
linguistic-resource representation~\citep{joshi2020state}; and World Bank income
groups~\citep{worldbank2025classification}.

The pre-specified sample of 15 (Brazil, Mexico, Argentina, Peru; Nigeria, South
Africa, Kenya, Egypt; India, Indonesia, Bangladesh, the Philippines; and the
United States, Germany, Japan as Global North references) covers four world
regions and four income groups. The linguistic-resource axis turned out far less
differentiated than intended: the dominant official languages fall into only
three Joshi classes, twelve of the fifteen sharing Class~5, and extending to 25
does not help, since eighteen are Class~5. This is a property of the world
rather than of our draw, because the official languages of most states are
themselves high-resource, but it leaves the linguistic axis too little variance
to be separated from the others, and we treat the Joshi results accordingly
(Section~\ref{sec:results:h1}). Countries with compromised statistical
governance or without a publicly retrievable national standard were excluded
from the pre-specified fifteen; Oceania entered only with the extension
(Australia), and Joshi Class~0 languages are not represented, which we note as a
scope limitation.

\subsection{Sample expansion after the pre-specified analysis}
\label{sec:expansion}

After the pre-specified 15-country analysis, and reported here as a post-hoc
extension rather than as part of the locked protocol, the sample was enlarged
to 25 countries under identical procedures, to raise power for the gradient and
mechanism tests and to multiply the within-country language pairs available for
H2. The manuscript reports the 25-country sample throughout; the primary family
recomputed on the pre-specified 15 alone is tabulated in Supplementary Material,
so that a reader can see what the locked protocol would have concluded. Because
the extension followed observation of the initial results, no pre-specified
conclusion is revised on the extension alone, and the expansion is treated as a
power-and-robustness check rather than as a new test. Angola, like Nigeria and
Argentina, has no national ambient PM\textsubscript{2.5} standard, adding
no-standard cases in which a faithful model must decline to state a non-existent
threshold. The ten additions and their roles are detailed in Supplementary
Material. Two extra question variants per (country, task) were generated for
the H2 pairs at design time but were never collected or scored; the H2 analysis
uses the single native-language rendering of each English prompt.

\subsection{Country covariates}
\label{sec:covariates}

Human Development Index (UNDP 2023--24) and GDP per capita index developmental
position and enter the mechanism analysis as adjustment covariates. The
corpus-representation exposures divide along a distinction the mechanism
analysis makes explicit, because the two channels apply to different effects.
Language-corpus measures quantify how much pretraining text exists in a language
and are the candidate mechanism for the native-language penalty: mC4 token count
(\citealp{xue2021mt5}, Table~6), OSCAR-2301 corpus size (corpus introduced by
\citealp{abadji2022oscar}; per-language figures read from the OSCAR-2301
dataset card, whose totals differ from the 22.01 release described in that
paper), and the Joshi class~\citep{joshi2020state}. Country-corpus measures
quantify how much the encyclopedic corpus is about a country, independent of
language, and are the candidate mechanism for the geographic gap: the English
Wikipedia article length, the number of Wikipedia language editions with an
article on the country, and the number of Wikidata statements on the country
entity.

The two families are not interchangeable. Because the geographic gap is
measured in English, with language held constant across countries, a residual
North/South gap cannot be a language-corpus effect. Wikipedia and Wikidata are
never ground truth here; they are the explanatory variable. What an answer is
scored against is always an official register (Section~\ref{sec:groundtruth}).
Had we scored against an encyclopedia, the study would measure agreement with an
encyclopedia, which is not the question a public administration faces.

\subsection{Model selection}

Fourteen LLMs were benchmarked across five access tiers, spanning roughly two
orders of magnitude in parameter count and four training-data geographies
(United States, China, Europe, Brazil): open-weight frontier (5), open-weight
mid (3) and open-weight small or regional (2), closed-accessible (3) and
closed-frontier (1). Three of the open-weight models were served locally
(Ollama, pinned tags: Phi-4~14B, Qwen3~14B, Cabra-Mistral~7B); the other seven
open-weight models were served through hosted endpoints (Groq free tier for
five, the DeepSeek and Cohere APIs for two), and the four closed models through
their vendors' APIs. The small tier includes a Brazilian-Portuguese
instruction-tuned model (Cabra-Mistral~7B~v3, a fine-tune of
Mistral-7B-Instruct-v0.3), tested against a scale-matched globally trained open
model for H3. Each model was queried with an identical pinned tag throughout the
collection window. The full roster is in Supplementary Table~S1.

We compare models as served through a fixed access stack, not idealised
weights. Provider infrastructure, serving configuration and API layer are part
of what each label denotes, and a residual gap could in principle reflect the
stack rather than the weights; we hold this constant by pinning version tags and
a fixed sampling configuration per provider throughout a fixed window
(Section~\ref{sec:collection}).

\subsection{Domain and task taxonomy}
\label{sec:tasks}

The single domain is air pollution policy, chosen because it anchors to
officially published, verifiable ground truth and maps to concrete decision
activities of a public environmental manager. Each of the five tasks corresponds
to a question a manager has to answer in the course of ordinary work, and each
was drafted so that a wrong answer would carry an administrative consequence.
Table~\ref{tab:tasks} gives them with the prompt exactly as issued to the models,
using Brazil as the worked instance, alongside the decision each stands in for.

\begin{table}[!t]
\centering
\caption{The five tasks, the prompt as issued, and the decision each stands in
for. Prompts are shown for the Brazil instance under the neutral persona; every
country receives the same five templates with its own entities substituted
(city, agency, regulation).}
\label{tab:tasks}
\footnotesize
\setlength{\tabcolsep}{5pt}
\renewcommand{\arraystretch}{1.0}
\begin{tabular}{@{}>{\raggedright\arraybackslash}p{0.11\linewidth}>{\raggedright\arraybackslash}p{0.54\linewidth}>{\raggedright\arraybackslash}p{0.28\linewidth}@{}}
\toprule
Task & Prompt as issued (Brazil) & Decision it stands in for \\
\midrule
T1 \newline Technical standard
& ``What is the current national annual mean ambient air quality standard for
PM\textsubscript{2.5} in Brazil? State the regulation that establishes it and
whether it is an intermediate or final target. Limit to 1--3 sentences.''
& Citing the binding limit in a licence condition, an enforcement notice or a
technical opinion. The wrong value voids the act. \\
\addlinespace
T2 \newline Local factual datum
& ``According to the national or subnational environmental agency, what was the
annual mean PM\textsubscript{2.5} concentration in São Paulo, Brazil, in the
most recent reported year? Provide value in
$\mu$g/m\textsuperscript{3}, the reference station or network, and the source.''
& Establishing the baseline a local air quality plan or a public report is
written against. \\
\addlinespace
T3 \newline Health-evidence synthesis
& ``Summarize the epidemiological evidence on mortality attributable to ambient
PM\textsubscript{2.5} exposure in Brazil over the last decade. Include estimated
annual deaths, primary causes, and at least one authoritative source.''
& Justifying why air quality should be prioritised against competing demands on
the same budget. \\
\addlinespace
T4 \newline Policy instruments
& ``List the principal national programs or legal instruments for monitoring and
controlling ambient air pollution in Brazil, with the executing agency for
each.''
& Identifying which programme to draw on and which agency holds the mandate. \\
\addlinespace
T5 \newline Applied recommendation
& ``A metropolitan region of Brazil faces a forecast episode with
PM\textsubscript{2.5} expected to exceed 75~$\mu$g/m\textsuperscript{3} within
24 hours. Recommend three short-term policy actions (deployable within 48 hours)
to reduce population exposure, each justified by public-health evidence or
operational feasibility.''
& Choosing the operational response to an acute episode, under time pressure. \\
\bottomrule
\end{tabular}
\end{table}

The tasks are ordered by how much of the answer is fixed by an official source.
T1, T2 and T4 admit verifiable ground truth and carry the primary accuracy
signal; T3 and T5 are rubric-scored open-generation tasks, where competent
answers legitimately differ. The ordering lets us ask whether models fail
uniformly or specifically where a published national value has to be recalled,
which is the question the results turn on. The country-by-task mapping of
official source classes and verifiability status is in the ground-truth
registry (Section~\ref{sec:groundtruth}), and the full 25-country prompt set is
released with the data.

\subsection{Persona manipulation}
\label{sec:persona}

Persona is a within-prompt factor with two levels, and the manipulation follows
the practice the vendors themselves document and recommend
(Section~\ref{sec:related}). The neutral condition presents the query as a plain
information request; the public\_manager\_env condition prepends a fixed role
frame identifying the requester as a municipal or regional environmental
management official, holding the substantive question identical. For Brazil's
T1 the frame is the single sentence \emph{``I am the municipal environmental
management official of São Paulo, Brazil, and I need to brief
decision-makers.''} A single template is applied uniformly across countries and
tasks, translated where the prompt language is non-English, so that any
between-condition difference is attributable to the frame rather than to
question wording. Persona order within each prompt-model cell is randomised,
and a manipulation check verifies that the frame is acknowledged in the
response. Exact wording, translation procedure and manipulation-check coding are
in Supplementary Material.

One implementation detail matters for how the result should be read. Every
request in this study is sent as a single user turn with no system prompt, so
the role frame is prefixed to the user message rather than installed as a
system instruction. Three of the four vendor guides place the role in the system
prompt; Google's states that it belongs either in the system instruction or
``at the very beginning of the user prompt''~\citep{google2026promptingstrategies},
which is the channel we test. The choice is deliberate: the system-prompt route
requires API access and a developer to configure it, while the public servants
whose decisions motivate this study reach these models through the consumer
chat interface, where the only available channel is the message they type. We
therefore test role framing as the sector performs it, and
Section~\ref{sec:limits} states what remains untested. Persona is the basis of
H6.

\subsection{Prompt construction and validation}

Prompts were authored for the domain by the research team with domain
supervision and validated against the ground-truth registry before entering the
main stimulus set, with every factual claim in a registry entry required to
resolve to an official source URL or a peer-reviewed DOI before the
corresponding prompt was admitted. The analysis plan called for a 50-prompt
pilot scored by two independent human raters ($\alpha\geq0.70$) before scaling
the design. That pilot was not run, and no rater-level $\alpha$ exists for this
study; rubric reliability is evidenced instead on the main data, by a
three-vendor panel over every item that required a judge
(Section~\ref{sec:results:reliability}), which returns $\alpha=\nalpha$. We
report this as a departure from the plan rather than as a threshold met.

All prompts were authored in English and, for a stratified subset of countries,
additionally in a relevant native language (Supplementary Section~S3).
Native-language renderings were produced by an LLM translator (Claude
Sonnet~4.6) prompted to preserve meaning, persona frame, proper nouns and
acronyms, with the factual ground truth held in English since the target value
is language-invariant. We did not employ human native-speaker translators, which
remains a limitation (Section~\ref{sec:limits}); we audited the renderings after
the fact by back-translating all 90 native prompts with a different model and
having two further models compare each against the English original
(Section~\ref{sec:results:h2}). The native language is the language of the
user, not necessarily of the instrument: Brazil's, Portugal's and the five
Spanish-speaking countries' standards are published in the prompt language,
whereas India's National Ambient Air Quality Standards are published in
English, so the Hindi condition asks in a language the official text does not
use.

\subsection{Ground truth and the ground-truth registry}
\label{sec:groundtruth}

Ground truth is a set of four task-specific registries, released with the data,
and we describe each as it was built rather than as it was planned
(Supplementary Table~S14 records the deviations). For T1 the key is the annual
PM\textsubscript{2.5} value in force during the collection window, taken from
the national instrument itself (statute, decree, regulation or gazette) and
stored with the source, a verbatim excerpt and a status flag. Twenty-one
countries have a numeric key. Three (Angola, Argentina, Nigeria) have no
national PM\textsubscript{2.5} standard, so the correct answer is to say so,
and a stated value is scored as fabricated. Egypt has no key, because the
official Arabic text could not be retrieved, and is excluded from T1 (56
cells). Bangladesh is scored against the 35~$\mu$g/m\textsuperscript{3} of its
Air Pollution (Control) Rules 2022, which replaced the 15 of the 2005 gazette;
an earlier version of the registry carried the 2005 value, and the correction
is logged with the released files. Where an instrument sets a ladder of staged
values (Brazil's five CONAMA stages, the EU limit of 25 with 10 from 2030, the
United Kingdom's 20 with an English target of 10 by 2040, the United States'
9.0 succeeding 12.0, Bangladesh's 35 succeeding 15), the registry stores the
whole ladder as well as the value in force, so that citing a neighbouring stage
can be separated from fabrication (Section~\ref{sec:results:t1}). Canada's
CAAQS and the Philippine value are objectives rather than legally binding
limits; they are the only national reference values those countries publish and
are treated as the key, which the text notes where it matters. The registry
does not record the date on which each instrument entered into force, so the
distance between that date and a model's training cut-off could not be entered
as a covariate; the ladder analysis in Section~\ref{sec:results:t1} is the
nearest available substitute.

For T2 the key is the most recent annual mean PM\textsubscript{2.5}
concentration that the WHO Ambient Air Quality Database (v6.1, 2024) records
for the prompted city, which compiles the values reported by national and
subnational agencies; any reported year is accepted within $\pm 20\%$, so that a
model citing an agency value from a different year is not penalised for the
year. The window is not tier-neutral by construction, because Global North
cities have longer WHO series (\nttwoyearsgn{} accepted years on average against
\nttwoyearsgs{}), but the range of values it admits is similar in the two tiers
(ratio of the largest to the smallest accepted value \nttwospreadgn{} against
\nttwospreadgs{}). Two cities have no PM\textsubscript{2.5} entry (Luanda,
Cairo) and those cells fall to the judge. For T3 the key is the WHO Global
Health Observatory estimate of deaths attributable to ambient air pollution
(indicator AIR\_41, 2021) with its uncertainty interval widened by half on each
side as the accepted range, plus the published ordering of causes; the
indicator covers ambient air pollution as a whole while the prompt asks about
PM\textsubscript{2.5}, and a model citing a different burden study (the GBD, a
national cohort) is scored by the judge on whether its figure and sources are
defensible, not against the WHO number. For T4 the reference set is the UNEP
Global Air Pollution Legislation catalogue for each country, with manual
overrides where it is stale. The registry records source, URL, access date and
status per entry, but it has no human-validator field and no per-document hash:
verification was documentary, by the authors, against the primary text. Entries
without an equivalent official source are excluded from the code-adjudicated
comparison for that cell, so that missing ground truth is never scored as a
model error.

\subsection{Data collection}
\label{sec:collection}

Calls were executed at temperature~0.3, with two replications per
prompt-model-persona combination, inside a fixed window to minimise model-update
confounds; the reported accuracies are therefore point-in-time estimates. Two
exceptions are declared. The GPT-5 family rejects any temperature other than the
default, so GPT-5 and GPT-5-mini ran at temperature~1.0 with reasoning effort
set to minimal; and no seed was sent to any provider, an earlier draft having
stated otherwise, so the two replicates at temperature~0.3 bound, rather than
remove, within-model variance (Supplementary Material compares the headline
effects with that variance). Every response was stored with complete request
metadata (pinned tag, timestamp, token counts, finish reason, persona condition)
under SHA-256 checksums. No response was excluded for its length or its
language: empty responses are scored as returning no value, and a response in a
language other than the one requested is scored on its content. Empty responses
are more frequent under native-language prompts (\nemptyen\% of English
responses on the three code-scored tasks against \nemptyes\% in Spanish,
\nemptypt\% in Portuguese and \nemptyhi\% in Hindi), which is part of the
language effect rather than attrition from it.

Per-cell coverage is not uniform: free and rate-limited endpoints returned
fewer admissible responses than metered APIs, so the number of scored
English-prompt cells ranges from 261 to 500 per model (777 for the model with
the most native-language pairs when those are counted). We report each model's
realised $N$ alongside its score (Table~\ref{tab:conf-model}) rather than
dropping under-covered cells, the tier gap is re-estimated on the prompts every
model answered (Section~\ref{sec:results:h1}), and the full coverage matrix,
with quota and exclusion reasons for every shortfall, is in Supplementary
Material.

\subsection{Measures}

Scoring assigns each task the most reliable instrument it admits.

Where the answer is a number in an official register, code decides. For T1, T2
and T3 the returned value is extracted and compared with the registry by a
deterministic scorer, with plausibility guards that reject implausible
extractions (a four-digit year read as a death count, a national population
read as a mortality figure, a 24-hour value read as an annual one). The
extractor normalises Devanagari digits, decimal commas and units written out in
words, and its behaviour on the cases that matter for the language comparison
(a comma decimal, a Hindi numeral, a parenthetical daily value beside the
annual one) is fixed by a unit-test suite released with the code; a stratified
manual audit of native-language T1 responses it classified as carrying no value
is reported in Supplementary Material, and no equivalent audit was run for T2
and T3. Comparing a value with a key is exact, whereas a language model asked to
do the same arithmetic reproduces it with noise. Code resolves \ncodecovtone\%
of the T1 cells in the analysis (\ntonecodecells{} of \ntonecells),
\ncodecovttwo\% of T2 and \ncodecovtthree\% of T3; a T1 response with no
extractable value counts as not having returned it. For the three no-standard
countries, code scores whether the answer states that no national standard
exists (correct), asserts a value (fabricated) or declines to answer (neither),
and the tier comparison is reported with and without them.

The code verdict replaces one subcomponent of the composite, factual accuracy
(weight $0.30$); the other four subcomponents of a T1, T2 or T3 cell remain
those assigned by the judge that scored it. For T1 that judge is the collection
judge, so the composite on T1 still carries the judge-and-tier confound
described below in $0.70$ of its weight. For that reason the primary outcome
for the tier comparison on the register-valued tasks is the binary code verdict
itself, whether the answer returned the value in force, and the composite is
reported as the secondary outcome (Section~\ref{sec:results:h1}).

Where judgement is required, a three-vendor panel decides by its mean. The
residual of T2 and T3, with all of T4, is scored by Gemini~2.5~Pro, Claude
Sonnet~4.6 and DeepSeek-V3, chosen for vendor diversity so that intra-vendor
correlation does not inflate apparent agreement. T5, the open recommendation,
retains the scores of the original collection judge; that judge was GPT-5-mini
for the fifteen pre-specified countries and Claude Haiku~4.5 for the ten added
in the extension, and because the extension holds seven of the ten Global North
countries, judge and tier are confounded in any effect that still rests on those
scores. This is why the factual verdict on T1 was moved to code: on the original
scores the Global South deficit in T1 was \njtgsgpt{} against \njtgngpt{} under
GPT-5-mini but \njtgshaiku{} against \njtgnhaiku{} under Claude Haiku, a
difference between judges as large as the difference between tiers
(Supplementary Material). T5 carries no tier effect in either instrument.
Reliability is reported as Krippendorff's $\alpha$, the single-judge ICC$(2,1)$
and, as the operative quantity, the ICC$(2,3)$ of the panel mean, a two-way
random-effects, absolute-agreement coefficient in which the judges' different
severities count against it, computed on every item the panel scored rather
than on a calibration subset. Language-model judges agree with human raters at
roughly the level humans agree with each other on open-ended
tasks~\citep{zheng2023judging}, but they carry position, verbosity and
self-enhancement biases, and models recognise their own outputs well enough to
prefer them~\citep{zheng2023judging,panickssery2024self}. One panel member
(DeepSeek-V3) is also an evaluated model, and we test for self-favouring
directly; the other two judges
belong to vendors whose models are in the roster (Anthropic, Google), so a
vendor-level preference cannot be excluded by that test. No human gold-standard
layer was collected; ground truth is anchored to official primary sources
(Section~\ref{sec:groundtruth}).

The unit of analysis is the scored cell, not the log line. Duplication exists at
two levels and each has a stated rule. At the response level, 21.9\% of (model,
prompt, replicate) keys have more than one stored answer, because collection
runs were resumed and re-run; the answer that enters the analysis is the first
stored in run-file order, which is the one the judge scored, and the code
verdicts are computed on that same answer so that the two instruments describe
one response. Scoring every stored answer instead and averaging the verdicts
changes the code verdict in \ndupchange\% of cells and the mean by \ndupmean
(Supplementary Material). At the score level, 951 cells (11.5\%) carry more
than one judge score, because the judge was re-run on parts of the set on
different dates; those are repeated measurements of the same item and are
collapsed by their mean, leaving 8{,}300 cells, or 8{,}244 once the 56 Egyptian
T1 cells without a key are removed (6{,}573 English-prompt). Treating re-runs
as independent would inflate $n$ and give extra weight to whichever cells
happened to be reprocessed. The rule was fixed on methodological grounds and
does not depend on the result.

The primary composite combines five pre-specified subcomponents (Supplementary
Section~S6), each normalised to $[0,1]$: factual accuracy against the registry
entry (weight $0.30$), contextual completeness ($0.25$), citation quality
($0.15$), calibration ($0.15$) and hallucination absence ($0.15$). Every
headline effect is recomputed under equal weights and a PCA-derived weighting,
and a conclusion counts as robust only if its direction is stable across all
three. The composite is computed per (country, model, persona) cell to support
the H6 difference-in-differences test.

\subsection{Analytical strategy}

The three primary tests are pre-specified as a single family (F1) and corrected
together with Bonferroni-Holm. Secondary tests (F2) are reported unadjusted
within family, and declared exploratory tests (F3) use Benjamini-Hochberg FDR
at $q=0.10$~\citep{benjamini1995fdr}, whose guarantee holds under positive
dependence; the Benjamini-Yekutieli correction, valid under arbitrary
dependence, returns the same decisions on this family, since the F3 tests it
covers have $p$-values several orders of magnitude below the threshold.

H1 (geographic gradient, primary) regresses country-level mean composite
accuracy on the Joshi/HDI gradient via Spearman $\rho$, pre-specified one-sided
at $\rho\geq0.55$, with a Mann-Kendall trend test as a mandatory
non-monotonicity complement. The binary Global North/South contrast on the same
country means was pre-specified as a sensitivity analysis of H1, and its
decomposition by task (Section~\ref{sec:results:h1}) was not pre-specified and
is exploratory. H4 (corpus mechanism, primary) relates corpus representation to
accuracy by partial Spearman $\rho$ adjusting for HDI and log GDP per capita,
with Wikipedia counts as the pre-specified proxy and Wikidata coverage as
exploratory complements. H6 (persona, primary) is a country-level
difference-in-differences, $\Delta_{\mathrm{GS}}-\Delta_{\mathrm{GN}}$,
supported if the persona condition reduces the Global South gap by at least
$0.05$ on the composite, with inference by a 5{,}000-permutation country-level
test; because the pre-specified test cannot distinguish a null from an
inconclusive result, we add an equivalence test, two one-sided tests on the
country-level bootstrap distribution of the difference-in-differences with
margins of $\pm2$ and $\pm5$~pp. H3 (secondary) is a one-sided Mann-Whitney
contrast (Cliff's $\delta$). H2 and H5 (exploratory) are reported descriptively
under F3 correction.

Wherever a test in the Results is first run on responses, we also report it
with the country or the model as the unit, because responses are nested in 25
countries and 14 models and a response-level $p$-value overstates the evidence
for any contrast that varies between countries. Three cross-cutting analyses
corroborate the tier gap: a Gaussian linear mixed model (\texttt{statsmodels}
MixedLM, REML) with crossed random intercepts for country and model, adjusting
for centred HDI; a Bayesian hierarchical re-estimation
(\texttt{pymc}~\citep{pymc}, weakly informative priors); and E-values for
robustness to unmeasured confounding, computed for the standardised mean
difference in composite accuracy between tiers through the approximation
$\mathrm{RR}\approx\exp(0.91\,d)$ that \citet{vanderweele2017evalue} prescribe
for standardised effects and whose assumptions it inherits. The corpus
mechanism is probed with an exploratory mediation model
(\texttt{semopy}~\citep{semopy}), and the persona manipulation with a
role-frame-acknowledgement check. Pre-specified robustness analyses appear in
Section~\ref{sec:results:robust}.

\subsection{Study protocol and ethics}

The analytic dataset, after quality gates but before model fitting, is frozen
and hashed, so that the analysis can be audited against the locked
specification rather than taken on trust. Release terms are stated at the end
of the paper.

No personal data from end-users of LLMs were collected; the study analyses
model outputs to publicly relevant policy questions. Provider terms of service
were reviewed and respected, and responses are redistributed solely for
research under fair-use and reproducibility provisions, excluded from any
commercial application.
